\documentclass[tikz,border=8pt]{standalone}
\usepackage{amsmath}
\usepackage{fontspec}
\usepackage{xeCJK}
\setmainfont{PT Serif}
\setCJKmainfont{Noto Sans CJK SC}
\usetikzlibrary{arrows.meta,positioning,calc}

\definecolor{blockblue}{HTML}{4A90D9}
\definecolor{sumgreen}{HTML}{5CB85C}
\definecolor{taporange}{HTML}{F0AD4E}
\definecolor{linegray}{HTML}{333333}

\tikzset{
  block/.style={rectangle, draw=blockblue, fill=none, line width=1.2pt,
    minimum height=0.75cm, minimum width=1.4cm, inner sep=2pt,
    font=\footnotesize\bfseries, align=center},
  wideblock/.style={block, minimum width=2.1cm},
  sum/.style={circle, draw=sumgreen, fill=none, line width=1.2pt,
    minimum size=0.34cm, inner sep=0pt,
    path picture={
      \draw[line width=0.8pt, linegray]
        (path picture bounding box.south west) -- (path picture bounding box.north east);
      \draw[line width=0.8pt, linegray]
        (path picture bounding box.north west) -- (path picture bounding box.south east);
    }},
  tap/.style={circle, draw=taporange, fill=none, line width=1.2pt,
    minimum size=4pt, inner sep=0pt},
  signal/.style={-{Latex[length=2.2mm]}, draw=linegray, line width=1.2pt},
  plain/.style={draw=linegray, line width=1.2pt},
  signal label above/.style={fill=none, above=4pt},
  neg sign/.style={fill=none, font=\footnotesize, very near end},
  note/.style={align=center}
}

\begin{document}
\begin{tikzpicture}[x=1cm,y=1cm]
  % 开环：增益只缩放输入输出关系，不进入对象分母。
  \node[note, font=\small\bfseries, align=right] (openlabel) at (0,1.55) {开环};
  \node (r1) at (1.0,1.55) {$r(t)$};
  \node[block] (k1) at (2.55,1.55) {$K$};
  \node[wideblock] (g1) at (4.95,1.55) {$G_0(s)=\dfrac{1}{s(s+4)}$};
  \node (y1) at (7.1,1.55) {$y(t)$};
  \node[note, font=\footnotesize, align=left] at (9.45,1.55)
    {开环传递函数：\\[2pt]
     $G_{\mathrm{OL}}(s)=\dfrac{K}{s(s+4)}$\\[2pt]
     极点固定为 $0,-4$};

  \draw[signal] (r1) -- (k1);
  \draw[signal] (k1) -- (g1);
  \draw[signal] (g1) -- (y1);

  % 闭环：单位负反馈使 K 进入闭环特征方程。
  \node[note, font=\small\bfseries, align=right] (closedlabel) at (0,-1.15) {闭环};
  \node (r2) at (1.0,-1.15) {$r(t)$};
  \node[sum] (sum) at (2.05,-1.15) {};
  \node[block] (k2) at (3.55,-1.15) {$K$};
  \node[wideblock] (g2) at (5.95,-1.15) {$G_0(s)=\dfrac{1}{s(s+4)}$};
  \node[tap] (tap) at (7.55,-1.15) {};
  \node (y2) at (8.5,-1.15) {$y(t)$};
  \node[note, font=\footnotesize, align=left] at (10.75,-1.15)
    {$\Phi(s)=\dfrac{K}{s^2+4s+K}$\\[2pt]
     特征方程：$s^2+4s+K=0$};

  \draw[signal] (r2) -- (sum);
  \draw[signal] (sum) -- node[signal label above, font=\footnotesize] {$e(t)$} (k2);
  \draw[signal] (k2) -- (g2);
  \draw[signal] (g2) -- (tap) -- (y2);
  \draw[plain] (tap) -- ++(0,-1.15) coordinate (returnright)
    -- coordinate[pos=0.52] (returnmid) ($(sum)+(0,-1.15)$);
  \draw[signal] ($(sum)+(0,-1.15)$) -- node[neg sign, left] {$-$} (sum.south);
  \node[note, font=\footnotesize, below=3pt of returnmid] {单位反馈 $H(s)=1$};
\end{tikzpicture}
\end{document}
